% Generated draft: requires mathematical and source review.
\documentclass[11pt]{article}
\usepackage[margin=25mm]{geometry}
\usepackage{fontspec}
\setmainfont{Latin Modern Roman}
\usepackage{amsmath,amssymb,mathtools,unicode-math}
\setmathfont{Latin Modern Math}
\usepackage{enumitem,needspace,xurl,hyperref}
\setlength{\parindent}{0pt}
\setlength{\parskip}{5pt}
\setlength{\emergencystretch}{3em}
\newcommand{\N}{\mathbb{N}}
\newcommand{\Z}{\mathbb{Z}}
\newcommand{\Q}{\mathbb{Q}}
\newcommand{\R}{\mathbb{R}}
\newcommand{\C}{\mathbb{C}}
\newcommand{\F}{\mathbb{F}}
\newcommand{\A}{\mathbb{A}}
\newcommand{\PP}{\mathbb{P}}
\newcommand{\E}{\mathbb{E}}
\newcommand{\eps}{\varepsilon}
\newcommand{\dd}{\,\mathrm{d}}
\begin{document}
\section*{UnsolvedMath ID: 30003326}
\textbf{Source problem code:} catalogue entry
\subsection{Definitions and mathematical statement}
Work in the language \(L_{\exp}=\{0,1,+,-,\cdot,<,\exp\}\) of ordered fields with a distinguished exponential function. Under the usual convention, an ordered exponential field is an ordered field \(K\) equipped with an order-preserving group isomorphism \(\exp_K:(K,+)\to(K_{>0},\cdot)\). The supplied record does not explicitly specify this convention or define its additional condition \(\mathrm{EXP}\); that condition must therefore remain an unresolved source qualification.

An \(L_{\exp}\)-structure is \emph{o-minimal} if every definable subset of its underlying ordered set, allowing parameters, is a finite union of points and intervals with endpoints in the structure or \(\pm\infty\). Two \(L_{\exp}\)-structures \(M,N\) are \emph{elementarily equivalent}, written \(M\equiv_{L_{\exp}}N\), if they satisfy the same \(L_{\exp}\)-sentences. An \emph{elementary embedding} \(j:M\to N\) is an embedding such that, for every \(L_{\exp}\)-formula \(\varphi\) and every finite tuple \(\bar a\) from \(M\),
\[
M\models\varphi(\bar a)\quad\Longleftrightarrow\quad N\models\varphi(j(\bar a)).
\]

A residue field associated with a valuation ring \(\mathcal O\subseteq K\), with maximal ideal \(\mathfrak m\), is \(\mathcal O/\mathfrak m\). If \(\mathcal O\) is convex, this quotient inherits an ordering. An induced exponential, when defined, has the form
\[
\exp_{\mathrm{res}}(a+\mathfrak m)=\exp_K(a)+\mathfrak m\qquad(a\in\mathcal O).
\]
The supplied record does not identify \(\mathcal O\), specify the intended residue construction, or give the hypotheses ensuring that this formula defines an exponential on the quotient. Thus \(K_{\mathrm{res}}\) below denotes the intended residue exponential field, whose precise construction remains to be recovered from the source.

\textbf{Question in the supplied current statement.} For every o-minimal ordered exponential field \(K\) satisfying the source condition \(\mathrm{EXP}\), is
\[
K_{\mathrm{res}}\equiv_{L_{\exp}}K\,?
\]
This is a conditional statement template, pending definitions of \(\mathrm{EXP}\) and \(K_{\mathrm{res}}\).

\textbf{Additional scope in the supplied original statement.} The original wording asks, for every such \(K\), whether all three assertions hold:
\[
\begin{gathered}
K_{\mathrm{res}}\text{ is o-minimal},\\
K_{\mathrm{res}}\equiv_{L_{\exp}}K,\\
\text{there exists an elementary embedding }j:K_{\mathrm{res}}\to K.
\end{gathered}
\]
The supplied current statement retains only the elementary-equivalence question. Both formulations are recorded here because the supplied evidence does not explain this change of scope. The language is also provisional if the source intends additional symbols.
\subsection{Short English statement}
The current question asks whether every o-minimal ordered exponential field satisfying the unspecified condition \(\mathrm{EXP}\) has the same first-order theory as its intended residue exponential field. The original wording additionally asks about o-minimality of the residue field and an elementary embedding into the original field. Definitions of \(\mathrm{EXP}\), the residue construction, and the intended language require clarification from the source.
\subsection{Sources}
\begin{itemize}
\item Local UnsolvedMath registry, supplied record with permanent ID \(30003326\), title \emph{Model Theory of Residue Exponential Fields}, collection \emph{Oberwolfach Reports Open Problems}, exact source problem code \texttt{catalogue entry}. The supplied record contains both the narrower current statement and the broader original statement. No individual proposer is supplied.
\item Supplied report citation: \emph{Mini-Workshop: Surreal Numbers, Surreal Analysis, Hahn Fields and Derivations (2017)}, Oberwolfach Reports, pp.\ 3313--3372, DOI \texttt{10.4171/owr/2016/60}. Supplied source URL and literature-source link, labeled ``original OWR report'': \url{https://doi.org/10.4171/owr/2016/60}. The record separately gives proposed year \(2016\) and tag \texttt{owr-year:2016}; their relationship to the citation's \(2017\) has not been resolved here.
\item Literature assessment dated exactly \texttt{2026-08-21}: this is supplied dated curation, not independently verified. It labels the record open and states that the extraction is only an open-questions heading without a standalone proposition, classified as open under a conservative rule. The background explicitly requests independent verification of time-sensitive status and cited claims.
\item The supplied statement-status label is \texttt{corrected\_verified}; its accompanying verification text reports that the official passage was checked and rewritten as a self-contained problem while preserving mathematical content. This is a supplied verification claim, not an independent verification in this draft, and it is in tension with the supplied heading-only literature assessment.
\end{itemize}
\end{document}
